\begin{lemma}\label{lem:multi-even-equality}
Suppose the host has positive even maximum degree $D$ and the
number of marked copies is $\frac{5D^2}{4}$. The host induced by
the marked endpoints is a blow-up of $C_5$ with all five parts of
size $\frac{D}{2}$. Every edge of that blow-up is single and marked.
\end{lemma}

\begin{proof}
The five-position multiplicity bound applies at $M=\frac{5}{4}$ and
gives $e<0.2105$. If every induced red five-cycle had
$\frac{R}{D}\ge\frac{1}{2}$, the support bound in \eqref{eq:support-projection-four} would give $S\le1.22181$ and
hence $e>0.0225$. These bounds satisfy the hypotheses of the last six-position certificate in \textbf{Table~\ref{tab:uniform-chain}}, whose upper bound is strictly less than
$\frac{5}{4}$. Thus there is a cycle with $R<\frac{D}{2}$.

Apply the redistribution leading to \eqref{eq:multi-rough-discharge}
at this cycle. For every outside vertex $x$ with $k_x>0$, its
resulting contribution is strictly less than $\frac{k_xD}{2}$.
For $k_x=1$, a vertex either keeps at most $R-2$ or sends half of
its marked degree and has a blue anchor edge. For $k_x=2$, there
is a blue anchor edge and no incoming transfer. For $k_x=3$ with
a marked neighbor of signature size zero, the incoming transfer
is at most $\max\{\frac{D+R}{4},R\}<\frac{D}{2}$. Without such a
neighbor, it is at most $\frac{D-3}{2}$, since the three anchor
adjacencies use degree units inside the core. For $k_x\ge4$,
the transfer is at most $D-k_x$ and its own contribution is
at most $D$. Vertices of signature size zero retain no positive
contribution. Since $K\le R$, the presence of an outside vertex
of positive signature size would give $2m<D|B|+\frac{DK}{2}
\le\frac{5D^2}{2}$ by \eqref{eq:multi-core-cut}.

Hence $K=0$. Now $2m\le D|B|=\frac{D(5D-R)}{2}$, so $R=0$.
By \eqref{eq:multi-core-identities}, the five part sizes are
$\frac{D}{2}$. For a vertex of $U_i$, the neighbors in
$U_{i-1}\cup U_{i+1}$ already use all $D$ degree units. There are
$\frac{5D^2}{4}$ edges on these pairs. Any marked copy outside
$B$ would require an outside vertex of positive signature size.
Thus all edges in the blow-up are marked, and $B$ contains every
marked endpoint.
\end{proof}


We finish with a consequence used in \textbf{Section~\ref{sec:exact}}
to force an induced red five-cycle. For $m,D>0$ we have
\begin{equation}\label{eq:absence-red-cycle}
 q_C=0 \quad\Longrightarrow\quad M<\frac{98}{81}.
\end{equation}
Suppose that $q_C=0$ and $M\geq\frac{98}{81}$. The uniform bounds for
the multigraph and its simple support give $M,S\leq\frac54$.
If $b=0$ then $e=0$. If $b>0$, \textbf{Lemma~\ref{lem:multi-extra}}
gives $e\leq1-\frac{M}{2}$. Thus in either case
\[
 \frac{98}{81}\leq M\leq\frac{12883}{10000},
 \qquad 0\leq e\leq\frac{32}{81}.
\]
These are the intervals used in the stored calculation.

Let $P$ be the five-position pattern of an induced red five-cycle.
As in the definition of $q_C$, every color-preserving map of $P$ is injective. Hence $q_C=0$ excludes every such map. It follows from
the definition of $X(P,T)$ in
\textbf{Subsection~\ref{sec:weighted-counts}} that $X(P,T)=0$ for
each forest $T\subseteq E(P)$ used in the calculation, including
$T=\varnothing$. Since the coordinates are nonnegative, the verifier
imposes all these equations by summing the distinct coordinates
$X(P,T)$ and requiring this sum to be at most zero.

Together with this inequality, the five-position calculation uses
the weighted degree and flag constraints, the simple-support bound,
the first inequality in \eqref{eq:support-projection-four}, and the
rooted five-cycle bounds indexed in the proof notes accompanying~\cite{FGSTSupplement}.
Using the five-position certificate from Section~G.4 of
the supplement~\cite{FGSTSupplement} and adding the correction for
negative residual coefficients gives
\[
 M\leq\frac{14802480264533669793521}{12250000000000000000000}
   <\frac{98}{81}.
\]
This contradicts the assumed range of $M$ and proves
\eqref{eq:absence-red-cycle}.
